\documentclass[12pt]{article} % Clase de documento: artículo y tamaño de letra
\usepackage[utf8]{inputenc} % Escritura en castellano con acentos
\usepackage[T1]{fontenc} % Escritura en castellano con acentos
\usepackage{calligra} 
\usepackage{pslatex}  % Fuente de letras
\usepackage{graphicx}
\usepackage[margin=2.5cm]{geometry}
\usepackage{txfonts}
\usepackage{xcolor}
\usepackage{fancyhdr}
\usepackage{graphicx}
\usepackage{float}
\usepackage{listings}
\usepackage{multicol}
\usepackage{color}
\usepackage{times}
\usepackage{tikz}
\usepackage{verbatim}
\usepackage{mdwlist}
\usepackage{array}
\usepackage{caption}

\usetikzlibrary{chains,fit,shapes,arrows,calc,shapes,decorations.pathreplacing}


\pagestyle{fancy}
\headheight=50pt %para cambiar el tamaño del encabezado
\fancyhead[L]    %la "L" indica a la izquierda
{	
 \begin{minipage}{3.3cm}
  \includegraphics[width=1.0\textwidth]{Logo-UTN-BA.jpeg}
%  \includegraphics[width=0.6\textwidth]{UTNlogo.eps}
 \end{minipage}	
 \begin{minipage}{7.7cm}
   \normalsize
   {
     \textsf
     {
       \calligra{Universidad Tecnológica Nacional\\ Facultad Regional Buenos Aires \\ Departamento de {Ingeniería} {Electrónica}} 
     }
   }
 \end{minipage}
}

\fancyhead[R] %la "R" indica a la derecha
{
  \begin{minipage}{4.0cm}
   \small
   {
     \emph{\textbf{Informática I}} \\ \emph{\today} \\ \emph{Ejemplos de parcial}\\ \emph{Curso r1092}  
   }
  \end{minipage}}


\begin{document} % Inicio del documento
\newpage
%Cambia de página, el texto después de este comando aparecerá en la siguiente página en adelante.


\noindent
  \begin{center}
   \begin{tabular}{| c | c | c |}
    \hline
     Apellido y Nombres \hspace{8cm} &  Legajo & {Calificación} \\ \hline 
      &	& \\ \hline
   \end{tabular}	
  \end{center}

\noindent
Todas las funciones pedidas deben ser colocadas en el archivo {\bf primerParcialFunciones.c} y debe estar acompañado por su correspondiente {\bf punto h}.\\
Además se debe escribir un programa de testeo para cada función requerida. El nombre de esos archivos es {\bf primerParcialNumeroEjercicio.c}.\\
\\
%%%%%%%%%%%%%%%%%%%%%%%%%%%%%%%%%%%%%%%%%%%%%%%%%%%%%%%%%%%%%%%%%%%%%%%%%%%%%%%%%%%%%%%%%%%%%%%%%%%%%%%%%%%%% Comienzo del tema del examen
%%%%%%%%%%%%%%%%%%%%%%%%%%%%%%%%%%%%%%%%%%%%%%%%%%%%%%%%%%%%%%%%%%%%%%%%%%%%%%%%%%%%%%%%%%%%%%%%%%%%%%%%%%%%
\lstset{
	frame=Ltb,
	framerule=0pt,
	aboveskip=0.5cm,
	framextopmargin=3pt,
	framexbottommargin=3pt,
	framexleftmargin=0.4cm,
	framesep=0pt,
	rulesep=.4pt,
	backgroundcolor=\color{gray97},
	rulesepcolor=\color{black},
% 	language=C,
	captionpos=b,
	tabsize=3,
	frame=lines,
	keywordstyle=\color{blue},
	commentstyle=\color{Gray},
	stringstyle=\color{red},
	numbers=left,
	numberstyle=\tiny,
	numbersep=5pt,
	breaklines=true,
	showstringspaces=false,
	basicstyle=\ttfamily\scriptsize,
	emph={label},
	framerule=0pt,
}

Implemente una función que obtenga el dígito verificador de un número de CUIT pasado como 
parámetro, el cálculo se realiza utilizando el algoritmo {\bf módulo11}.\\
 El prototipo de la función es el siguiente:
\begin{center}
{\bf int cuit\_validar (char *cuit);}
\end{center}
El parámetro {\bf cuit} es un puntero al vector que contiene el número de CUIT terminado en ‘{\bf $\setminus$0}’\\
\\
Devuelve:
\begin{itemize}
\item Un {\bf número positivo} indicando el dígito verificador. 
\item {\bf -1}: cuando la cantidad de dígitos es distinto de 10 
\item {\bf -2}: Indica que el número de CUIT es inválido (contiene algo distinto a números) \end{itemize}
Algoritmo módulo 11:  
\begin{itemize}
\item Multiplique los dígitos Desde el menos significativo por la serie 2,3,4,5,6,7.
\item Sume el resultado de las multiplicaciones anteriores.
\item Calcule el módulo 11 de la suma anterior.
\item Al resultado anterior restele 11
\begin{itemize}
\item si el resultado es {\bf menor que 10} lo obtenido es el dígito verificador
\item si el resultado es {\bf 10} el dígito verificador es {\bf 9}.
\item Si el resultado es {\bf 11} el dígito verificador es {\bf 0}.
\end{itemize} 
\end{itemize}
Ejemplo:\\
\\
%\vspace{0.3cm}
\noindent
\fontsize{10pt}{10pt}\selectfont
   \begin{tabular}{| c | c | c | c | c | c | c | c | c | c | c | c | c | c |}
    \hline
     CUIT & 2 & 0 & 1 & 2 & 3 & 4 & 5 & 6 & 7 & 8 & suma & \%11 & dígito\\ \hline 
     Valor a multiplicar por digito & $X_5$ & $X_4$ & $X_3$ & $X_2$ & $X_7$ & $X_6$ & $X_5$ & $X_4$ & $X_3$ & $X_2$ &  &148\%11 & 11-5 \\ \hline 
     Resultado de la multiplicación & 10 & 0 & 3 & 4 & 21 & 24 & 25 & 24 & 21 & 16 & =148 & =5 & 6\\ \hline 
   \end{tabular}	
\fontsize{12pt}{12pt}\selectfont
 
\end{document} % Fin del documento.
 

  
